\chapter{Open Problems}
\label{ch:open}

The book has proposed a vocabulary and proved what follows from it. It has also, repeatedly, marked claims as principles and left questions to experiment. This chapter collects the principal open problems, sorted by the kind of work needed to settle them. Each is stated so that a result would be recognizable.

\section{Mathematical problems}

\emph{A sheaf-theoretic account of editing.} Chapter~\ref{ch:sheaves} treated the sheaf of takes over a time interval, and Chapter~\ref{ch:descent} treated coverage over a scene. A unified treatment would place both on a single site whose objects are spacetime regions, with a sheaf of scene states and a morphism of sheaves into images. The open problem is to characterize the image of rendering in the category of sheaves in terms of cohomological conditions, so that a statement of the form ``this family of images lifts to a consistent world'' has a cohomological criterion.

\emph{Higher obstructions.} The class in $H^1$ of Chapter~\ref{ch:cohomology} detects the failure of a single correction to be explained by local data. Obstructions to the existence of descent data in the stronger sense of Chapter~\ref{ch:descent} live in higher cohomology or in nonabelian cohomology with coefficients in groups of isomorphisms of the local objects. Whether such higher obstructions have natural filmic interpretations, beyond identity-tracking across several views, is unsettled.

\emph{Typed generation.} The Galois connection of Chapter~\ref{ch:scripts} and the typing of Chapter~\ref{ch:types} describe satisfiability. The efficient generation of films from a rich type system, with decidable checking and tractable sampling, is a problem at the boundary of programming languages and probabilistic inference. For the linear fragment of Chapter~\ref{ch:linear}, the existence of balanced runs is a Petri-net reachability problem, and the cost of sampling uniformly from balanced runs is not addressed here.

\emph{Quantitative preservation.} \Cref{prop:compose-contracts} gives a worst-case bound on composite drift. Under distributional assumptions on the viewer's dissimilarity, one expects tighter bounds, and a theory of preservation for stochastic transformations would give them.

\section{Empirical problems}

\emph{Conditional dimension.} Principle~\ref{pr:conditional} says that the number of perceptually independent variations is a conditional quantity. The experiment is to estimate the singular spectrum of the sensitivity matrix for several scenes, tasks, and viewer groups, and to ask whether the spectra are similar enough across scenes to justify any common figure. A rapidly decaying spectrum for most scenes would support the cue-repertoire conjecture of Chapter~\ref{ch:equivalence}, and a flat spectrum would count against it.

\emph{The geodesic principle.} Principle~\ref{pr:geodesic} asserts that natural camera movements are critical points of an action with a perceptual metric. The test is to fit the metric to viewer judgments of naturalness for a set of synthetic camera paths and to ask whether the fitted metric makes held-out judgments predictable. A failure of predictive power for every metric in the Fisher--Rao family would count against the principle.

\emph{Recovery and occupation.} The temporal organization of acoustic occupation, formalized in Chapter~\ref{ch:contracts}, supports the hypothesis that arrangements with equal level histograms but different block structure differ in listener fatigue and comprehension. Whether they do, and whether any benefit of recovery intervals depends on their length and placement, are questions for controlled listening studies with the response variables of Chapter~\ref{ch:equivalence} specified in advance.

\emph{Authority of the history.} Principle~\ref{pr:authority} states that in a persistent generative world the history must be authoritative at commit. A systematic comparison of generative systems with and without an explicit history layer, on tasks requiring long-range consistency, would test whether the failures attributed to its absence are in fact explained by it.

\section{Design problems}

\emph{Interfaces for hierarchical direction.} Chapter~\ref{ch:direction} showed that prefix-free command sets parse uniquely and that commands generally do not commute. The design of notations that expose conflicts, preview realizations at the user's level of abstraction, and support undo without silently invalidating later commitments is an open problem in human-computer interaction.

\emph{Presentation contracts.} The same film may be offered under different presentation contracts, each preserving different things: theatrical form, quiet home viewing, analytical study. How to let a viewer move between them without losing place or narrative context is a problem about synchronization of several versions, which is a cover of the film by its presentations and, in the language of Chapter~\ref{ch:selective}, a problem of branching with shared segments.

\section{Limits of the approach}

The formalism treats a film as a path, a word, or a family of local data, and a viewer as a source of dissimilarity judgments. These are idealizations, and they exclude much of what criticism attends to. Interpretation is contestable even when the measurements are reproducible. A theorem about a model of a film is a theorem about the model, and the book's principles are the places where it ventures a claim about film itself, together with a statement of what would count against the claim. If the programme is to be useful, it will be because the models sharpen disagreements into questions that evidence can answer, and because they make clear what a generative system must be told if it is to avoid errors that no amount of rendering quality can repair.
